\section{Algolean}\label{Algolean}
Algolean\cite{algolean} ist ein Framework für Algorithmen und Komplexitätstheorie
und sein Kern ist die Freie Monade\cite{the_free_monad} Prog.
Der Name Algolean ist ein Wortspiel und kann auf zwei Arten gelesen werde.
Entweder als \glqq Algo\grqq ~+ \glqq lean\grqq, also im Sinn einer Algorithmus Bibliothek in Lean, oder \glqq Algol\grqq ~+ \glqq ean\grqq, als Hommage an Algol, wobei dann \glqq ean\grqq ~so ausgesprochen wird, dass sich der Name als Algolsch übersetzen lässt.

Die Grundidee des Frameworks ist es,
den Nutzer bei der Formalisierung seiner Algorithmen zu zwingen, seine Grundoperationen und deren Kosten im Vorhinein zu bestimmen.
Die Komplexität beziehungsweise Laufzeit der Algorithmen werden dann als direkte strukturelle Konsequenz der Grundkosten abgeleitet.

Der Grund, warum dieses Framework überhaupt benötigt wird ist, dass man nicht einfach eine Funktion definieren kann,
die eine Funktion und einen Input übergeben bekommt, und dann einfach zählt, wie oft bestimmte Funktion beim Ausführen aufgerufen werden.
Denn Lean sieht zwischen intentional verschiedenen Funktionen, die also für verschiedene Zwecke gedacht sind und beim Ausführen komplett unterschiedliche Dinge tun,
keinen Unterschied, solange sie extentional gleich sind, also für denselben Input denselben Output liefern.
Es wären also nicht erkennbar, welche Funktion genau übergeben wurde.

Aus diesem Grund arbeitet Algolean mit der Freien Monade.
Monaden sind in funktionalen Programmiersprachen der Weg um Computerberechnungen mit Seiteneffekten zu modellieren.
Das prominenteste Beispiel für eine Monade ist IO, über die zum Beispiel alle Interaktionen mit dem Standard Input und Output laufen.
% \cite{monad_wiki}

Einer der Seiteneffekte der Freien Monade ist, dass sie beim Aufrufen von Funktionen speichert, dass sie diesen Aufruf getätigt hat.

Innerhalb von Monaden wird die do-Notation benutzt.
Ähnlich dem Taktikmodus, der mit by geöffnet wird, beginnt diese mit \textbf{do}.
Die do-Notation kann das sequenzielle Abarbeiten sowie For- und Whileschleifen einer imperativen Programmiersprache simulieren.
Für die do-Notation und für Monaden typisch ist außerdem, das Zuweisen mit $\leftarrow$, was mit den Typen zusammenhängt und auf das ich später noch einmal zu sprechen komme.
Dabei ist unter der Haube alles funktionale Programmierung.

Mehr Informationen zu Monaden im Algemeinen finden sich in dem Buch Functional Programming in Lean\cite{functional_programming_in_lean} von David Thrane Christiansen
und zur Freien Monade im Speziellen im Blogbeitrag The Free Monad\cite{the_free_monad} von Tanner Duve.

Was außerdem sehr für dieses Framework spricht, ist wie einfach man bestimmen kann, was genau man zählen möchte.
Man möchte in der theoretischen Informatik ja nicht immer die bloße Anzahl der aufrufe einer Funktion zählen.
Je nachdem, was einen interessiert, kann man beispielsweise arithmetische Operationen ignorieren, als ein Aufruf zählen oder einen Wert anhand ihrer bitweisen Komplexität zuweisen.

Die für meine Arbeit am relevantesten Elemente von Algolean erkläre ich im Folgenden anhand des Beispiels listLinearSearch aus dem Algolean Repository von Shreyas Srinivas und Eric Wieser,
welches ich etwas angepasst habe, damit es anschaulicher wird.
Es zeigt die Laufzeitanalyse für die lineare Suche.
\begin{lstlisting}
def List.search {α : Type} [BEq α] (l : List α) (x : α) : Bool :=
  match l with
  | [] => false
  | head :: tail =>
    if head == x
    then true
    else search tail x
\end{lstlisting}

Die grundlegenden Funktionen, die man zählen möchte, werden in einer \textbf{Query} deklariert.
Eine Query ist eine induktive Funktion von Type nach Type.
Jede grundlegende Funktion bekommt dabei einen eigenen Konstruktor und hat als Typ die Query angewendet auf den Rückgabetyp der Funktion.
Im Beispiel hat die Query, ListSearch $\alpha$ nur einen Konstruktor, nämlich den, für das Vergleichen von einem Element a und einer Liste l und da der Rückgabetyp eines Vergleich Bool ist, ist der Typ des Konstruktors ListSearch $\alpha$ von Bool.
\begin{lstlisting}
inductive ListSearch (α : Type) : Type → Type _ where
  | compare (x y : α) : ListSearch α Bool
\end{lstlisting}

Ein Modell (engl. \textbf{Model}) einer Query definiert dann, wie sich die Funktionen verhalten und was die jeweiligen Kosten sind.
Dabei ist es möglich mehrere verschiedene Modelle für dieselbe Query zu definieren.
Der Typ eines Modells für eine Query Q und einen Kostentyp c ist dann Model Q c.
Im Beispiel vergleicht die Funktion, ob zwei Elemente gleich sind und die Kosten für einen solchen Vergleich sind konstant 1.
\begin{lstlisting}
def ListSearch.natCost {α : Type} [BEq α] : Model (ListSearch α) ℕ where
  evalQuery
    | compare x y => x == y
  cost _ := 1
\end{lstlisting}

Ein Algorithmus / Programm für eine bestimmte Query Q mit einem Rückgabetyp $\alpha$ wird jetzt als Monade vom Typ \textbf{Prog} Q $\alpha$ in der do-Notation definiert.
Dabei können dann die in der Query definierten Funktionen verwendet werden, wobei man darauf achten muss, dass man $\leftarrow$ benutzt, um einer Variablen den Rückgabewert der Funktion zuzuweisen.
Tut man das nicht und verwendet stattdessen :=, dann hat das Ergebnis den falschen Typ, nämlich den, den man beim Definieren der Query angegeben hat.
Im Beispiel wäre das dann ListSearch $\alpha$ Bool anstelle von Bool.
Was auch zu beachten ist, ist, dass man den Algorithmus für eine Query und nicht für ein Modell schreibt.
Das bedeutet, dass die Funktionen deklariert aber nicht definiert sind.
Dementsprechend muss man, wenn man genauere Informationen haben möchte, wie das Ergebnis zustande kommt, einen komplexeren Rückgabetyp definieren, in dem die Informationen mitgegeben werden können.
Worauf es beim Schreiben des Algorithmus dann ankommt, ist, dass man auch die entsprechenden Queryaufrufe verwendet und keine äquivalenten normalen Funktionen,
denn nur die Queryaufrufe werden am Ende gezählt.
Wenn man also keine Queryaufrufe verwendet, wird man als Ergebnis bekommen, dass der Algorithmus eine Laufzeit von 0 hat.
Im Beispiel ist der Algorithmus die lineare Suche.
Er iteriert also über die Liste und prüft in jedem Schritt mit der Vergleichsfunktion der Query, ob das aktuelle Element dem gesuchten entspricht.
\begin{lstlisting}
def listLinearSearch {α : Type} (l : List α) (x : α) : Prog (ListSearch α) Bool := do
  match l with
  | [] => return false
  | head :: tail =>
    let cmp : Bool ← ListSearch.compare head x
    if cmp then
      return true
    else
      listLinearSearch tail x
\end{lstlisting}

Wenn man einen Algorithmus jetzt ausführen will, oder seine Laufzeit berechnen, dann kann man sich nun ein zur Query passendes Modell wählen und den Algorithmus für dieses Modell auswerten.
Dafür gibt es die Funktion \textbf{Prog.eval} und \textbf{Prog.time}.
Da der ganze Aufbau mit Query, Modell und Monade einen großen Overhead bedeuten, ist die Intention bei Algolean,
dass man zu dem Algorithmus als Prog den Algorithmus noch ein weiteres Mal als normale Lean Funktion implementiert,
und anschließend zeigt, dass für das entsprechende Modell die beiden äquivalent sind.
Dadurch ist es möglich, die ganzen Beweise zur Korrektheit ohne diesen Overhead zu führen.
Im Beispiel wird als erstes gezeigt, dass die als Prog geschriebene lineare Suche äquivalent zur Funktion List.search ist.
Die Aussage hier per Induktion über die Funktion List.search gezeigt.
Dabei gibt es zwei Basisfälle und einen Induktionsschritt.
Im ersten Basisfall ist die Liste leer, was in beiden Funktionen false bewirkt.
Im zweiten Basisfall ist der Kopf der Liste das gesuchte Element.
Die compare Funktion gibt nach der Definition im Modell also true aus und damit auch die Funktion listLinearSearch.
Das passt also auch zu List.search.
Im Induktionsschritt sind der Kopf der Liste und das gesuchte Element verschieden.
Es rufen sich also beide Funktionen rekursiv auf.
Damit entspricht das Ziel genau der Induktionshypothese.
\begin{lstlisting}
lemma listLinearSearch_eval {α : Type} [BEq α] (l : List α) (x : α) :
    (listLinearSearch l x).eval ListSearch.natCost = l.search x := by
  fun_induction l.search x with
  | case1 =>
    simp [listLinearSearch]
  | case2 head tail h =>
    simp [listLinearSearch, ListSearch.natCost]
    simp [h]
  | case3 head tail h ih =>
    unfold listLinearSearch
    simp [ListSearch.natCost] at ⊢ ih
    simp [h]
    rw [← ih]
\end{lstlisting}

Anschließend wird gezeigt, dass die Laufzeit maximal der Länge der Liste entspricht, ebenfalls per Induktion.
Im ersten Fall ist die Liste leer und die Laufzeit 0, da compare nie aufgerufen wird.
Im zweiten Fall wird das Element gefunden.
Es wird also ein Vergleich angestellt, aber die Liste enthält auch mindestens ein Element.
Im dritten Fall, also dem Induktionsschritt wird das Element nicht gefunden.
Nach dem Vergleich wird also weiter gesucht.
In dem Fall ist die Laufzeit dann 1 + die Zeit für die Suche über die Restliste, aber die Länge der Liste ist auch 1 + die Länge der Restliste.
Wenn man also auf beiden Seiten 1 abzieht, entspricht das der Induktionshypothese.
\begin{lstlisting}
lemma listLinearSearchM_time_complexity_upper_bound {α : Type} [BEq α] (l : List α) (x : α) :
    (listLinearSearch l x).time ListSearch.natCost ≤ l.length := by
  fun_induction l.search x with
  | case1 =>
    simp [listLinearSearch]
  | case2 head tail h =>
    simp [listLinearSearch, ListSearch.natCost]
    simp [h]
  | case3 head tail h ih =>
    simp [listLinearSearch, ListSearch.natCost, h] at ⊢ ih
    lia
\end{lstlisting}

Ein weiteres Konzept ist die Reduktion (\textbf{Reduction}).
Diese bietet, für zwei Querys Q$_1$ und Q$_2$, die Möglichkeit Programme, die für Query Q$_1$ geschrieben wurden, als Programme für Q$_2$ zu betrachten,
indem sie jede Funktion aus Q$_1$ als ein für Q$_2$ geschriebenes Programm definiert.
Hierfür gibt es allerdings noch kein Beispiel, da Reduktionen im Framework zwar schon definiert sind, aber nach meinem Kenntnisstand, ich der erste bin, der sie in einem Projekt verwendet.

Mehr Informationen zu den verschiedenen Aspekten von Algolean gibt es in der dazugehörigen Dokumentation\cite{algolean_doc}.
